\subsection*{Backward angular-moment certification from the A.11 endpoint}
\label{subsec:a11-backward-angular-endpoint-20261021}
The source's original stress expression on the corrected azimuthal hull is
\[
Q_s=-1+(1-h)\frac{I}{XH}-D\eta\frac{I_\eta}{XH},
\qquad D=\tfrac12-h.
\]
The last term is essential within a correction support: the moment
and correction coefficients need not be independent of \(\eta\).

On the full closed A.11 hull, of logarithmic length \(23/10\), set
\(E=E_b(1+u)\), \(u=\sum_i c_i(\eta)\beta_i(y)\).
The unedited profile is independent of \(\eta\) there, with
\(\ell_b=-\lambda\). The two smooth-step-derivative bumps are
disjoint and satisfy
\(\|\beta_i\|_\infty\le640/3\) and
\(\|\beta_i'\|_\infty\le1715200/9\).
Assume the authentic root-jet inputs already derived for
\(T_f=512\), \(0<\lambda\le1/100\):
\[
|c_i|\le54\lambda^{28},\qquad
|c_{i,\eta}|\le6\lambda^{28},\qquad
0<h\le\lambda/100.
\]
Let \(y_R\) be the right edge of the second support. The
source's exact A.11 correction makes \(I/(XH)=(1-\lambda)^{-1}\)
at the end of the corrected plateau. From \(y_R\) to that
endpoint there is no edit and \(\ell=-\lambda\). Uniqueness for
\(r_y+(1-\lambda)r=1\), \(r=I/(XH)\), therefore gives the
*backward terminal data*
\[
r(y_R,\eta)=\frac1{1-\lambda},\qquad
r_\eta(y_R,\eta)=0.
\]
This is stronger than merely bounding the entrance moment.

\begin{theorem}[Whole-hull \(Q_s\) without a pointwise mesh]
\label{thm:a11-backward-angular-20261021}
Under the preceding authentic coefficient-jet and terminal-moment
hypotheses, the following hold throughout both correction supports
and the \(17/10\) intervening gap:
\[
\left|r-\frac1{1-\lambda}\right|<6\cdot10^8\lambda^{28},
\quad |r_\eta|<13\cdot10^7\lambda^{28},
\quad \left|\frac{I_\eta}{XH}\right|<131\cdot10^6\lambda^{28},
\]
\[
-\frac{3\lambda}{2}<\ell<-\frac\lambda2,\qquad
Q_s\ge\frac{\lambda-h}{1-\lambda}
-7\cdot10^8\lambda^{28}>\frac{49}{50}\lambda.
\]
\end{theorem}
\begin{proof}
Disjointness and the source bounds imply
\[
|u|\le11520\lambda^{28}<1/100,\quad
|u_y|\le10291200\lambda^{28},\quad
|u_\eta|\le1280\lambda^{28},\quad
|u_{y\eta}|\le(10291200/9)\lambda^{28}.
\]
Because \(\ell+\lambda=u_y/(1+u)\),
\(g:=\partial_\eta\log E=u_\eta/(1+u)\), and
\(\ell_\eta=u_{y\eta}/(1+u)-u_yu_\eta/(1+u)^2\),
exact rational comparisons at \(\lambda=1/100\) give
\[
|\ell+\lambda|<11\cdot10^6\lambda^{28},\quad
|g|<1300\lambda^{28},\quad
|\ell_\eta|<13\cdot10^5\lambda^{28}.
\]
The moment ratio satisfies \(r_y+(1+\ell)r=1\).
Writing \(r_*=(1-\lambda)^{-1}\), integrate this equation
backward from \(r(y_R)=r_*\):
\[
r(y)-r_*
=r_*\int_y^{y_R}\exp\left(\int_y^t(1+\ell(v))\,dv\right)
(\ell(t)+\lambda)\,dt.
\]
The hull length \(23/10\) and \(e^3<21\) imply
\[
|r-r_*|<\frac{100}{99}\frac{23}{10}
21\cdot11\cdot10^6\lambda^{28}
<6\cdot10^8\lambda^{28}.
\]
In particular \(0<r<2\). Differentiating the transport
equation and integrating backward from \(r_\eta(y_R)=0\)
gives
\[
|r_\eta|<2\frac{23}{10}21\cdot13\cdot10^5\lambda^{28}
<13\cdot10^7\lambda^{28}.
\]
Since \(I_\eta/(XH)=r_\eta+rg\), its bound follows.
Substitution into the *full* source formula for \(Q_s\)
gives an error below \(7\cdot10^8\lambda^{28}\).
Finally \(700000000(1/100)^{27}<1/100\), and
\(h\le\lambda/100\). The slope interval follows from
the first derivative bound.
\end{proof}

\begin{corollary}[Strict cone from a separate \(N_s/E\) tail gate]
\label{cor:a11-backward-cone-20261021}
In addition assume *for this same corrected profile* that
the exact \(S(\infty)=\Pi(\infty)=0\) gauges and authenticated
future-integral bounds give \(|N_s|/E\le503\lambda/2000\),
and \(X\ge9/\lambda\) on the hull. Then, with
\(a=2-2\ell\), \(p_1=XQ_s/(1-2h\eta^2)\) and
\(w=N_s/(EQ_s)\),
\[
|N_s|/E<\lambda/2,\qquad
2(1-a/p_1)^2-(a-2)w^2\ge
\frac{85676973}{76832000}>\frac{11}{10}>\frac9{20}.
\]
\end{corollary}
\begin{proof}
The theorem gives \(p_1\ge441/50\),
\(a\le43/20\), \(a/p_1<1/4\),
\(|w|\le503/1960\), and \(a-2\le3/20\).
The displayed value is the exact rational consequence.
\end{proof}

The advance removes the need to separately interval-sample
\(Q_s\) or the corrected slope on the A.11 hull,
*conditional on the same authentic root and terminal moment
gauge*. It does **not** authenticate the A.19 \(S\)-root,
the terminal pressure gauge, the future \(N_s\)-tail,
the full exterior stage covering, Appendix B's nonlinear
contraction, the unforced startup hierarchy, or Theorem~1.1.
The companion Python checker verifies exact rational consequences;
the new Lean module contains only scalar gates and does not
formalize these source integral/ODE/PDE arguments.